\section{Syntax and semantics}

\begin{figure}[p]
  \begin{alignat*}{2}
         &\text{Types } \tau, \tau_0, \tau_1 &::=~ &\mathbf{unit}
         \mid \num[i, f]
         \mid \tau_0 \times \tau_1
         \mid \tau_0 \otimes \tau_1
         \mid \tau_0 + \tau_1
         \mid \tau_0 \multimap \tau_1
         \mid {\bang{s} \tau}
         \mid {M_u \tau}
         \\
         &\text{Values } v, w \ &::=~ &\langle \rangle
         \mid k \in \textit{num}
         \mid \langle x,y \rangle 
         \mid (x, y)
         \mid \tin_i \ x
         \mid \lambda x.~e
         % \mid \mathbf{op}(x) \\
         % & & & \mid [v]
         % \mid\rnd v
         % \mid {\ret v} 
         % \mid \factor v
         % \mid \letbind x = v \ \tin \ f 
         \\
         &\text{Terms } e, f, g &::=~ &v
         \mid x
         % \mid (b_0, b_1)
         \mid \mathbf{op} ~ e
         \mid e~f
         \mid {\pi}_i\ e
         \mid \langle e,f \rangle 
         \mid (e, f)
         \mid \tin_i \ e \\
         & & & \mid \letpair \ (x,y) = e \ \tin \ f
         \mid \letassign x  = e \ \tin \ f 
         \mid [e] \\
         & & & 
         \mid 
          \mathbf{case} \ e \ \mathbf{of} \ (\tin_1 \ x.f_1 \ | \ \tin_2 \ x.f_2)
         \mid \rnd e
         \mid {\ret e} 
         \mid \factor e \\
         & & & 
         \mid {\letbind x = e \ \tin \ f}
         \mid \letcobind x = e \ \tin \ f
  \end{alignat*}
  \caption{
    Types, values, and terms. $\mathbf{op} \in \mathcal{O}$. $i \in \{1, 2\}$.
    $\textbf{op}$ is a syntactic metavar representing
    functions of arbitrary arity.
    $\textit{num}$ is a parameter in the programming language representing the
    set of numbers the language computes over.
  }
  \label{fig:syntax}
\end{figure}

\begin{figure}[p]
\begin{center}
%% unit
\AXC{}
\RightLabel{(Unit)}
\UIC{$\Gamma \vdash \langle \rangle : \mathbf{unit}$}
\bottomAlignProof
\DisplayProof
\hskip 0.5em
%% var
\AXC{$s \ge 1$}
\RightLabel{(Var)}
  \UIC{$\Gamma, x:_s \tau, \Theta \vdash x : \tau$}
\bottomAlignProof
\DisplayProof
\hskip 0.5em
%% fun
\AXC{$\Gamma, x:_1 \tau_0 \vdash e : \tau$}
\RightLabel{($\multimap$ I)}
\UnaryInfC{$\Gamma \vdash \lambda x. e : \tau_0 \multimap \tau $}
\bottomAlignProof
\DisplayProof
\vskip 1em

%% app
\AXC{$\ \Gamma \vdash e : \tau_0 \multimap \tau$}
\AXC{$\ \Theta \vdash f : \tau_0 $}
\RightLabel{($\multimap$ E)}
\BinaryInfC{$\ \Gamma + \Theta \vdash ef : \tau $}
\bottomAlignProof
\DisplayProof
\hskip 0.5em
%% dep prod intro
\AXC{$\Gamma \vdash e : \tau_0$}
\AXC{$\Gamma \vdash f : \tau_1$}
\RightLabel{($\times$ I)}
\BinaryInfC{$\Gamma \vdash \langle e, f \rangle: \tau_0 \times \tau_1$}
\bottomAlignProof
\DisplayProof
\vskip 1em
%%


%% ROW2
%% dep prod elim
\AXC{$\Gamma \vdash e : \tau_1 \times \tau_2$}
\RightLabel{($\times$ E)}
\UIC{$\Gamma \vdash {\pi}_i \ e : \tau_i$}
\bottomAlignProof
\DisplayProof
\hskip 0.5em
%% ind prod intro
\AXC{$\Gamma \vdash e : \tau_0 $}
\AXC{$\Theta \vdash f : \tau_1$}
\RightLabel{($\tensor$ I)}
\BIC{$\Gamma + \Theta \vdash (e, f) : \tau_0 \tensor \tau_1$}
\bottomAlignProof
\DisplayProof
\vskip 1em

%% ind prod elim
\AXC{$\Gamma \vdash e : \tau_0 \tensor \tau_1$ }
\AXC{$\Theta,x:_s \tau_0,y:_s\tau_1 \vdash f: \tau $}
\RightLabel{($\tensor$ E)}
\BIC{$s * \Gamma + \Theta \vdash \letpair (x,y) \ = \ e \ \tin \ f : \tau $}
\bottomAlignProof
\DisplayProof
\hskip 0.5em
%% ind sum intro
\AXC{$\Gamma \vdash e : \tau_0$ }
\RightLabel{($+$ $\text{I}_i$)}
\UIC{$\Gamma \vdash \mathbf{in}_i \ e : \tau_0 + \tau_1$}
\bottomAlignProof
\DisplayProof
\vskip 1em

% sum elim
\AXC{$\Gamma \vdash e : \tau_0+\tau_1$}
\AXC{$\Theta, x:_s \tau_0 \vdash f_1 : \tau$ \qquad
$\Delta \ | \ \Theta, x:_s \tau_1 \vdash f_2: \tau$}
\RightLabel{($+$ E)}
\AXC{$s > 0$}
\TIC{$s * \Gamma + \Theta \vdash \mathbf{case} \ e \ \mathbf{of} \ (\mathbf{in}_1 x.f_1 \ | \ \mathbf{in}_2 x.f_2) : \tau$}
\bottomAlignProof
\DisplayProof
\vskip 1em

% box elim
\AXC{$\Gamma \vdash e : {!_s \tau_0}$}
\AXC{$\Theta, x:_{t*s} \tau_0 \vdash f : \tau$}
\RightLabel{($!$ E)}
\BIC{$t * \Gamma + \Theta \vdash \letcobind x = e \ \tin \ f : \tau$}
\bottomAlignProof
\DisplayProof
\hskip 0.5em
% ops
\AXC{$\{ \mathbf{op} : \tau_0 \multimap \tau_1 \} \in \Sigma$}
\RightLabel{(Op)}
\UIC{$\Gamma \vdash \mathbf{op} : \tau_0 \multimap \tau_1 $}
\bottomAlignProof
\DisplayProof
\vskip 1em
%%


%% ROW 5

%%% ROW 6

% let 
\AXC{$\Gamma \vdash e :  \tau_0$}
\AXC{$\Theta, x:_{s} \tau \vdash f : \tau$}
\RightLabel{(Let)}
\BIC{$s * \Gamma + \Theta \vdash \letassign x = e \ \tin \ f : \tau$}
\bottomAlignProof
\DisplayProof
\hskip 0.5em
%% const
% TODO: need to describe inj function to some degree here, I suppose
\AXC{$k \in \textit{num}$}
\RightLabel{(Const)}
\UIC{$\Gamma \vdash k : \num[i, f]$}
\bottomAlignProof
\DisplayProof
\vskip 1em

% box intro
\AXC{$\Gamma \vdash e : \tau$ }
\RightLabel{($!$ I)}
\UIC{$s * \Gamma \vdash [e] : {!_s \tau}$}
\bottomAlignProof
\DisplayProof
\vskip 1em

%%% ROW 7

%% subsumption
\AXC{$\Gamma \vdash e :  M_q \tau$}
\AXC{$r \ge q$}
\RightLabel{(Subsumption)}
\BIC{$\Gamma \vdash e :  M_{r} \tau$}
\bottomAlignProof
\DisplayProof
\hskip 0.5em
% factor
\AXC{$\Gamma \vdash e : (M_q \tau_0) \times (M_r \tau_1)$}
\RightLabel{(Factor)}
\UIC{$\Gamma \vdash \factor \ e : M_{max(q,r)} (\tau_0 \times \tau_1)$}
\bottomAlignProof
\DisplayProof
\vskip 1em
%% return
\AXC{$\Gamma \vdash e : \tau$}
\RightLabel{(Ret)}
\UIC{$\Gamma \vdash \ret e : M_0 \tau$}
\bottomAlignProof
\DisplayProof
\hskip 0.5em
%% RND
\AXC{$\Gamma \vdash e : \num[i,f+1]$}
\RightLabel{(Rnd)}
\UIC{$\Gamma \vdash \rnd \ e : M_u \ \num[i,f]$}
\bottomAlignProof
\DisplayProof
\vskip 1em


%%% ROW 8


% let-bind
\AXC{$\Gamma \vdash e : M_r \tau_0$}
\AXC{$\Theta, x:_{s} \tau_0 \vdash f : M_{q} \tau$}
\RightLabel{($M_u$ E)}
\BIC{$s * \Gamma + \Theta \vdash \letbind x = e \ \tin \ f : M_{s*r+q} \tau$}
\bottomAlignProof
\DisplayProof
\vskip 1em

\vskip 5em

\end{center}
    \caption{Typing rules for \Lang, with $s,t,q,r,u \in \NNR \cup \{\infty\}$
      and for $i \in \{ 1, 2 \}$ where $u$ is a fixed constant parameter (see
      Definition~\ref{def:modular-interface} for details on picking an adequate
      constant).}
    \label{fig:typing_rules}
\end{figure}

\begin{figure}[p]
  \begin{alignat*}{2}
         &\text{Evaluation contexts } C &::=~ &[.] 
         \mid \mathbf{op}(C) 
         \mid C~e 
         \mid v~C 
         \mid \pi_i~C 
         \mid \langle C, e \rangle
         \mid \mathbf{in}_i~C \\
         & & &
         \mid \langle v, C \rangle
         \mid (C, e)
         \mid (v, C)
         \mid [ C ]
         \mid \mathbf{rnd}~C 
         \mid \mathbf{ret}~C 
         \mid \mathbf{factor}~C \\ 
         & & &
         \mid \letassign x = C \ \tin \ f
         \mid \letpair x = C \ \tin \ f \\
         & & &
         \mid \letbind x = C \ \tin \ f
         \mid \letcobind x = C \ \tin \ f \\
  \end{alignat*}
  \caption{Grammar for evaluation contexts.}
  \label{fig:eval-contexts}
\end{figure}

\begin{figure}
\begin{center}
\begin{equation*}
\begin{aligned}[c]
  \mathbf{op} \ v &\mapsto op \ v \\
	\pi_i\langle v_1,v_2 \rangle &\mapsto v_i \\
	(\lambda x.e) \ v &\mapsto e[v/x] \\
  \rnd k &\mapsto (k, \rho(k))
\end{aligned}
\quad
\begin{aligned}[c]
	\letassign x = v \ \tin \ e &\mapsto e[v/x] \\
  \letpair (x, y) = (v, w) \ \tin \ e &\mapsto e[v/x][w/y] \\
  \letcobind x = [v] \ \tin \ e &\mapsto e[v/x] \\
  \ret v &\mapsto (v, v)
\end{aligned}
\end{equation*}
\vskip -1em
\begin{align*}
	\mathbf{case} \ (\mathbf{in}_i \ v) \ \mathbf{of} \ (\mathbf{in}_1 \ x.e_1 \ | \ \mathbf{in}_2 \ x.e_2 )  &\mapsto e_i[v/x]
   \\
  \factor ((v_1, v_2), (v_3, v_4)) &\mapsto ((v_1, v_3), (v_2, v_4)) \\
  \\
  \Lambda \epsilon . e~{c} &\mapsto e[c /\epsilon] \\
\end{align*}
  \vskip -0.75em
  \AXC{$f[v_1/x] \mapsto^* (v_3, v_4)$}
  \AXC{$f[v_2/x] \mapsto^* (v_5, v_6)$}
  \BIC{$\letbind x = (v_1, v_2) \ \tin \ f \mapsto (v_3, v_6)$}
  \DisplayProof
	
  \vskip 0.4em
  \RightLabel{(Context Rule)}
  \AXC{$e \mapsto e'$}
  \UIC{$C[e] \mapsto C[e']$}
	\DisplayProof
\end{center}
    \caption{Substitution-style evaluation rules for \Lang. Parameterized for $i
    \in \{1, 2 \}$. $op$ is a higher-order metavar. When bolded \textbf{op}
    refers to the syntax and when italicized $op$ refers to the corresponding
    function on syntactic values (it may be a constant function).}
    \label{fig:sub_eval_rules}
\end{figure}
